\documentclass[11pt]{article}
\usepackage[margin=1in]{geometry}
\usepackage{xcolor}
\usepackage{enumitem}
\usepackage{hyperref}
\usepackage{booktabs}
\usepackage{amsmath}

\newcommand{\reviewer}[1]{\subsection*{#1}}
\newcommand{\point}[1]{\paragraph{#1}}
\newcommand{\response}{\textcolor{blue}{\textbf{Response:}}~}

\title{Author Response: ``When Does Conformer Geometry Help? Selective Complementarity of 3D Ensemble Statistics and 2D Fingerprints''}
\author{Anonymous Authors}
\date{}

\begin{document}
\maketitle

We thank all three reviewers for their thoughtful and constructive feedback. Below we address each point in detail. All changes in the revised manuscript are marked in \textcolor{blue}{blue text}.

%======================================================================
\reviewer{Reviewer 1P2E (Rating: 4, Borderline Accept)}
%======================================================================

\point{R1.1: Supplement empirical findings with a theoretical discussion explaining why simple covariance invariants outperform more complex representations (bias-variance trade-off, parameterization complexity, dataset size).}

\response We have substantially expanded the ``simplicity principle'' discussion in Section~6. The revised text now provides an explicit parameter-count analysis: the 5 scalar invariants operate at ${\sim}180$ samples per covariance parameter on ESOL ($N_\text{train} \approx 900$), whereas the lowrank variant ($k{=}128$) operates at ${\sim}7$---well below the threshold for reliable estimation. We connect this to classical model selection theory, showing that the bias-variance tradeoff penalizes higher-dimensional representations at the sample sizes typical of molecular property datasets (642--4,200 molecules). The new feature dimension ablation (Table~6, Section~5.6) provides independent confirmation: even for the mean $\boldsymbol{\mu}$, RMSE monotonically degrades from $D{=}256$ to $D{=}2{,}048$, demonstrating that this is a fundamental statistical constraint rather than an engineering limitation.

\point{R1.2: Conduct supplementary experiments on classification benchmarks.}

\response We have added classification experiments on BACE (1,513 molecules, $\beta$-secretase inhibition) and BBBP (2,038 molecules, blood-brain barrier permeability) in a new Section~5.5. Results (Table~7) confirm the property taxonomy extends to classification: conformer features hurt BACE by $-4.4\%$ AUROC and show no benefit on BBBP ($+0.1\%$). This is consistent with the taxonomy---$\beta$-secretase binding is an electronic/binding property, and BBB permeability is governed primarily by topological features. Per-seed results are provided in Appendix~K.

\point{R1.3: Validate key findings on larger or industrially sourced datasets.}

\response We acknowledge this as an important limitation in the revised Conclusion. Our benchmarks span 642 to 133K molecules across academic datasets (MoleculeNet, QM9, MARCEL); validation on larger industrially sourced datasets (ChEMBL, ZINC, proprietary ADMET panels) would strengthen the taxonomy's generalizability. We note this as a priority for future work and have added explicit language to this effect in Section~7 (Limitations).

\point{R1.4: Increase figure resolution, standardize metrics, correct typographical errors.}

\response We have conducted a full proofreading pass. All figures are exported at $\geq$300 DPI. Evaluation metrics are standardized (RMSE throughout for regression, AUROC for classification). Typographical errors and formatting inconsistencies have been corrected across tables and main text (see blue-marked changes throughout).

\point{R1.5: Provide clearer descriptions of scaffold split generation and number of replicates.}

\response We have expanded Section~3.7 (Training Details) with explicit descriptions of the Murcko scaffold split procedure (Bemis--Murcko generic scaffolds extracted via RDKit, sorted by frequency, assigned to splits maintaining 80/10/10 ratios) and seed counts (3 seeds for the main benchmark, 10 seeds for statistical validation experiments).

%======================================================================
\reviewer{Reviewer tcvY (Rating: 3, Borderline Reject)}
%======================================================================

\point{R2.1: Insufficient 3D GNN baselines---only SchNet is used; need DimeNet++, PaiNN, Equiformer.}

\response We have added PaiNN~[Sch\"utt et al., 2021] as a second 3D GNN baseline (Section~5.4, Table~5, Appendix~L). PaiNN uses equivariant message passing with separate scalar/vector channels---a fundamentally different inductive bias from SchNet's continuous-filter convolutions. Results confirm the hierarchy: PaiNN achieves 28\% (ESOL), 42\% (FreeSolv), and 30\% (Lipophilicity) gains over FP+XGBoost from a \emph{single} conformer, while SchNet uses 10.

The convergence of 21--42\% gains across these two architecturally diverse GNNs provides strong evidence that the hierarchy reflects a genuine property of the representation, not an artifact of SchNet specifically. We explain in the revised Section~5.4 why DimeNet++ (intermediate between SchNet and PaiNN in design) and Equiformer (primarily beneficial in large-scale pre-training settings orthogonal to our study) are not included: two architectures spanning the continuous-filter and equivariant paradigms are sufficient to establish robustness.

Regarding the reviewer's concern that SchNet could be trivially extended with ensemble pooling: this is precisely what our DKO mean ensemble baseline does (mean pooling over conformers through a neural architecture). The fact that even this approach underperforms FP+XGBoost strengthens, rather than undermines, the claim that the pre-computed feature bottleneck is the limiting factor.

\point{R2.2: FreeSolv results are directionally inconsistent ($-13.5\%$ to $+18.0\%$).}

\response We have clarified this in the revised Section~5.4. The $+18.0\%$ degradation comes from a 3-seed random-split auxiliary experiment on $N{=}642$ molecules---a regime where a single unfavorable split can dominate the mean. The definitive result is the 10-seed paired one-sided $t$-test on the main scaffold-split pipeline: $-13.5\%$ improvement, $p < 3 \times 10^{-5}$ (Table~4). The learning curve analysis (Section~5.7) further clarifies: hybrid features hurt at small data fractions but become beneficial above ${\sim}440$ molecules, explaining the sensitivity to split and seed variation at FreeSolv's small sample size. The 10-seed result, not the 3-seed auxiliary, should be taken as definitive.

\point{R2.3: Feature extraction pipeline (zero-padding/truncation to fixed $D$) may be the primary bottleneck.}

\response We have added a new discussion paragraph in Section~5.6 addressing this concern directly. The feature dimension ablation (Table~6) demonstrates that RMSE \emph{monotonically degrades} from $D{=}256$ to $D{=}2{,}048$. If truncation were the bottleneck, we would expect larger $D$ to help---instead, it hurts. This shows the limitation is the curse of dimensionality (noise dimensions that shallow models cannot exploit), not information loss from truncation. We acknowledge that BDE's variable dimensions (187--1,854) are a genuine concern and note adaptive-length encoding as future work, while emphasizing that this would not address the fundamental finding that pre-computed summary statistics discard the relational structure preserved by end-to-end GNNs.

\point{R2.4: Excluding pre-trained models limits generalizability of the taxonomy.}

\response We agree that pre-trained 3D models (e.g., Uni-Mol) may alter the taxonomy boundaries. We have expanded the Related Work (Section~2) and Limitations (Section~7) to explicitly scope the taxonomy to the non-pre-trained setting and acknowledge this limitation. Importantly, our DKO conformer statistics are \emph{complementary} to pre-trained representations---they could be concatenated with Uni-Mol embeddings just as they are concatenated with fingerprints here. Testing this combination is an important direction for future work, but is orthogonal to our central question of when 3D conformer geometry \emph{per se} adds predictive value beyond 2D fingerprints.

\point{R2.5: Restriction to regression tasks limits the framework's claimed generality.}

\response We have addressed this with new classification experiments on BACE and BBBP (Section~5.5, Table~7). Conformer features hurt BACE ($-4.4\%$ AUROC) and show no benefit on BBBP ($+0.1\%$), consistent with the taxonomy (binding and permeability are not solvation-dependent). We have updated the scope statement in the Introduction (Section~1) and abstract to explicitly include classification targets. We note that neural DKO variants require architectural adaptation for classification (different loss function, output activation) and are evaluated only on regression; the hybrid FP+conformer approach is evaluated on both.

%======================================================================
\reviewer{Reviewer u34v (Rating: 4, Borderline Accept)}
%======================================================================

\point{R3.1: Manuscript feels too fragmented due to titled paragraphs.}

\response We have revised the Discussion (Section~6) to reduce fragmentation. Bold titled paragraphs (``An empirical property taxonomy,'' ``Relationship to prior work,'' ``Practical implications'') have been converted to flowing prose with topic sentences that integrate the content more naturally. The section now reads as a cohesive narrative rather than a series of independent blocks.

\point{R3.2: Research questions formulated as actual questions feels too informal.}

\response We have rephrased RQ1--RQ3 in the Introduction as declarative objectives rather than interrogative questions (e.g., ``We characterize which molecular property types benefit\ldots'' rather than ``Which molecular property types benefit\ldots?'').

\point{R3.3: Performance hierarchy is confused and could benefit from clearer presentation.}

\response We have substantially rewritten the performance hierarchy presentation in Section~6. The revised text walks through each of the four tiers with explicit evidence, explains the progression from tier 4 to tier 1 in terms of how much 3D relational structure is preserved (summary statistics $\to$ physical invariants $\to$ full atomic graph), and provides the mechanistic insight that moving up the hierarchy corresponds to richer preservation of geometric information. The key result---that FP+3D descriptors match SchNet on ESOL (1.000 vs.\ 1.004)---is highlighted as evidence that carefully engineered features can approach end-to-end learning.

\point{R3.4: Some text too direct; expand into more technical discourse.}

\response We have expanded terse statements throughout the manuscript into more technical language, particularly in the abstract (performance hierarchy), Discussion (bias-variance analysis, hierarchy explanation), and Conclusion (decision framework). These changes are marked in blue.

\vspace{1em}
\noindent\rule{\textwidth}{0.4pt}
\vspace{0.5em}

\noindent\textbf{Summary of changes:}
\begin{itemize}[itemsep=2pt]
    \item Added PaiNN as second 3D GNN baseline (Section~5.4, Table~5, Appendix~L)
    \item Added classification benchmark on BACE and BBBP (Section~5.5, Table~7, Appendix~K)
    \item Added feature dimension ablation $D \in \{256, 512, 1024, 2048\}$ (Section~5.6, Table~6)
    \item Expanded bias-variance theoretical analysis (Section~6)
    \item Clarified FreeSolv directional inconsistency (Section~5.4)
    \item Addressed featurization bottleneck concern (Section~5.6)
    \item Expanded pre-training scoping and limitations (Sections~2, 7)
    \item Rewrote Discussion for clarity and flow (Section~6)
    \item Rephrased research questions as declarative objectives (Section~1)
    \item Clarified scaffold split procedure and replicate counts (Section~3.7)
    \item Added larger-dataset limitation acknowledgment (Section~7)
    \item All changes marked in \textcolor{blue}{blue text}
\end{itemize}

\end{document}
